\documentclass[a4paper,12pt]{article}
\usepackage{graphicx}
\usepackage{caption}
\usepackage{subcaption}
\usepackage{float}
\usepackage[margin=1in]{geometry}
\usepackage{lmodern}
\usepackage{amsmath}
\usepackage{xcolor}
\usepackage{listings}
\usepackage{pgffor}

\lstset{
  breaklines=true,
  postbreak=\mbox{\textcolor{red}{$\hookrightarrow$}\space},
}

\graphicspath{{../../kNN/output/}}

\begin{document}

\title{Relatório kNN}
\author{Eduardo}
\date{\today}
\maketitle

\section{Metodologia}
O código foi escrito em matlab 2024b.

\subsection{Algoritmo de classificação}
Foi implementado o algorimo \texttt{k-nearest neighbors} ponderado com \textit{kernel} arbitrário, conforme \ref{lst:kNN}.

\begin{lstlisting}[language=Matlab, caption=Algoritmo kNN, label=lst:kNN]
function label = mykNN(x, complete_set, k, h)
  % mykNN: binary classifier: label is {-1, 1}
  % x: point to be classified
  % complete_set: training set with class labels
  % k: number of nearest neighbors
  % h: kernel bandwidth
  % label: class label of x

  % number of training samples
  training_set_point_qtty = size(complete_set, 1);

  % number of features
  feature_qtty = size(complete_set, 2) - 1;

  % initialize distances vector
  distances = zeros(training_set_point_qtty, 1);

  % compute distances
  for i = 1:training_set_point_qtty
    distances(i) = euclidean_distance(x, complete_set(i, 1:feature_qtty));
  end

  % sort distances
  [sorted_distances, sorted_indexes] = sort(distances);

  % initialize class labels vector
  class_labels = zeros(k, 1);

  % get class labels of k nearest neighbors
  for i = 1:k
    class_labels(i) = complete_set(sorted_indexes(i), feature_qtty + 1);
  end

  % compute kernel weights
  weights = zeros(k, 1);

  for i = 1:k
    weights(i) = kernel_function(sorted_distances(i), h);
  end

  % compute class label
  label = sign(sum(weights .* class_labels));

end

function distance = euclidean_distance(x, y)
  % euclidean_distance: computes the euclidean distance between two points
  % x, y: points
  % distance: euclidean distance between x and y

  distance = sqrt(sum((x - y) .^ 2));

end

function value = kernel_function(distance, h)
  % kernel_function: computes the kernel function value
  % distance: distance between two points
  % h: kernel bandwidth
  % value: kernel function value

  value = exp(-distance / (2 * h ^ 2));

end

\end{lstlisting}

\subsection{Geração de dados}
Para avaliar o classificador, dados são gerados pelas funções em \ref{lst:data_generation}.

\begin{lstlisting}[language=Matlab, caption=Geração de dados, label=lst:data_generation]
function [data1, labels1, data2, labels2] = generate_dataset(datasetType, n, noise)
  % Generate synthetic dataset using MATLAB's built-in functions
  % datasetType: Choose between 'blobs', 'spirals', 'moons', 'xor', 'circles'
  % n: Number of samples
  % noise: Noise level

  switch datasetType
    case 'blobs'
      data1 = mvnrnd([2, 2], [0.5, 0; 0, 0.5], n/2) + noise * randn(n/2, 2);
      data2 = mvnrnd([-2, -2], [0.5, 0; 0, 0.5], n/2) + noise * randn(n/2, 2);
    case 'spirals'
      t = linspace(0, 4*pi, n/2)';
      data1 = [t.*cos(t), t.*sin(t)] + noise * randn(n/2, 2);
      data2 = [-t.*cos(t), -t.*sin(t)] + noise * randn(n/2, 2);
    case 'moons'
      [data, labels] = make_moons(n, noise);
      data1 = data(labels == 1, :);
      data2 = data(labels == 0, :);
    case 'xor'
      data1 = [randn(n/4, 2) + 2; randn(n/4, 2) - 2] + noise * randn(n/2, 2);
      data2 = [randn(n/4, 2) + [-2, 2]; randn(n/4, 2) - [-2, 2]] + noise * randn(n/2, 2);
    case 'circles'
      [data, labels] = make_circles(n, noise);
      data1 = data(labels == 1, :);
      data2 = data(labels == 0, :);
    otherwise
      error('Unknown dataset type');
  end

  % Assign labels
  labels1 = ones(size(data1, 1), 1);
  labels2 = -ones(size(data2, 1), 1);
end

% Helper functions for generating datasets
function [data, labels] = make_moons(n, noise)
  theta = linspace(0, pi, n/2)';
  r = 1 + noise * randn(n/2, 1);
  data1 = [r .* cos(theta), r .* sin(theta)];
  data2 = [1 - r .* cos(theta), -r .* sin(theta) - 0.5];
  data = [data1; data2] + noise * randn(n, 2);
  labels = [ones(n/2, 1); zeros(n/2, 1)];
end

function [data, labels] = make_circles(n, noise)
  theta = linspace(0, 2*pi, n/2)';
  r1 = 1 + noise * randn(n/2, 1);
  r2 = 0.5 + noise * randn(n/2, 1);
  data1 = [r1 .* cos(theta), r1 .* sin(theta)];
  data2 = [r2 .* cos(theta), r2 .* sin(theta)];
  data = [data1; data2] + noise * randn(n, 2);
  labels = [ones(n/2, 1); zeros(n/2, 1)];
end

\end{lstlisting}

\subsection{Calculo numérico da superfície de decisão}
Para permitir a visualização da superfície de decisão, o classificador, treinado, classifica diversos pontos ao longo da região de interesse. \ref{lst:grid_generation} produz a malha de pontos.

\begin{lstlisting}[language=Matlab, caption=Geração de malha de pontos, label=lst:grid_generation]
function [X, X1, X2] = generate_grid(complete_set, resolution)

  % Determine the range of the data
  min_x1 = min(complete_set(:, 1));
  max_x1 = max(complete_set(:, 1));
  min_x2 = min(complete_set(:, 2));
  max_x2 = max(complete_set(:, 2));

  % Generate a grid of points within the data range
  x1 = linspace(min_x1 - 1, max_x1 + 1, resolution);
  x2 = linspace(min_x2 - 1, max_x2 + 1, resolution);
  [X1, X2] = meshgrid(x1, x2);
  X = [X1(:), X2(:), zeros(length(X1(:)), 1)];

end

\end{lstlisting}

\section{Testes e Resultados}
O classificador foi avaliado para diferentes pares de \texttt{k} e \texttt{r} - a abertura da gaussiana geradora dos pontos de treinamento. Para cada par, a superfície de decisão é calculada, plotada e salva como imagem. \ref{lst:main} mostra o código principal.

\begin{lstlisting}[language=Matlab, caption=Código principal, label=lst:main]
% Generate synthetic dataset using MATLAB's built-in functions
datasetType = 'blobs'; % Choose between 'blobs', 'spirals', 'moons', 'xor', 'circles'

% Dataset parameters
samples = 100;
noise = 2;

% kNN parameters
k = 5;
h = 1;

for k = 1:20:50
  for noise = 1:1:3

    % Generate the dataset
    [data1, labels1, data2, labels2] = generate_dataset(datasetType, samples, noise);
    complete_set = [data1, labels1; data2, labels2];

    % Generate a grid for plotting
    [X, X1, X2] = generate_grid(complete_set, 100);

    % Classify the points
    for i = 1:size(X, 1)
      X(i, 3) = mykNN(X(i, 1:2), complete_set, k, h);
    end

    % Plot the decision boundaries
    plot_boundaries(X, X1, X2, data1, data2, true, string(k) + '_' + string(noise) + '.png');

  end
end

\end{lstlisting}

As figuras a seguir relacionam as superfícies de decisão aos parametros avaliados.
  
\foreach \k in {1, 21, 41} {
  \foreach \r in {1, 2, 3} {
    \begin{figure}[H]
      \centering
      \includegraphics[width=0.5\textwidth]{\k_\r.png}
      \caption{Superfície de decisão para k=\k\ e r=\r}
    \end{figure}
  }
}

\section{Conclusão}
Pelos resultados, fica claro que, quanto maior a quantidade de vizinhos avaliada, maior a robustês do classificador em relação a \textit{outliers}. No entanto, valores muito altos podem fazer com que este não capture com fidelidade a distribuição espacial dos dados.

\end{document}